\scanpage{34}

de façon précise

\[
\begin{aligned}
\widetilde C_\infty/\mu_\infty&\simeq\text{Image inverse de }\overbrace{\overline C_1/\mu_3}^{\substack{\omega\\\unsure{\wr|\simeq}}}\subset\mu_3^3/\mu_3=\Pi_3\\
&\hspace{5em}\text{par}\quad\Pi_\infty\mathrel{\unsure{-|\longrightarrow}}\Pi_3
\end{aligned}\tag*{(30)}
\]

En \unsure{vertu|outre|termes} de la suite exacte

\[
0\mathrel{-}\Pi_\infty\mathrel{\overset{3}{\longrightarrow}}\Pi_\infty
\mathrel{\overset{p_3}{\rule{1.5em}{0.4pt}}}
\underset{\substack{\cup\\\omega}}{\Pi_3}\mathrel{-}0
\tag*{(31)}
\]

$\Pi_\infty$ opère librement sur cette \unsure{``}partie de $\Pi_\infty$, image inverse de $\omega$ par $\Pi_\infty\mathrel{-}\Pi_3$. Ceci donne les opérations de $\Pi_\infty$ sur $E(\unsure{\sigma_*|O_*|O_{\sigma}|\alpha_*})\deleted{\unsure{\{a\}}}\simeq p_3^{-1}(\omega)$, celles de $\Gamma$ et $\unsure{\mathfrak S_3|\mathscr S_3}$ \deleted{étant} induites par leurs opérations sur $\Pi_\infty$ ....

\uline{NB} Sur $\Pi_3$ $\deleted{\unsure{E\times|\mathcal E\times}}\mathfrak S_3\times\Gamma$ opère via $\mathfrak S_3\times\overbrace{\{\pm1\}}^{\substack{\text{opère par}\\\text{homothéties}\\\text{sur }\Pi_3}}$, ($\Gamma\mathrel{\overset{\chi_3}{\rule{1.5em}{0.4pt}}}\{\pm1\}$). Considérons $\mathfrak S_3^+\inserted{\simeq\mathbb Z/3\mathbb Z}$ opérant sur $\Pi_3$ ($\simeq\mathbb F_3^2$), alors on sait qu'il doit y avoir une droite invariante -- \deleted{et les} en fait elle est unique -- et $\omega$ est formé des deux éléments non nuls de cette droite ....

% archive sheet number: 33
